\documentclass[twocolumn]{aastex631}
\begin{document}
\title{The reionization ionizing-photon-budget shortfall for star-forming galaxies is not robust to systematics at z\~{}7}
\author{NebulaMind Lab (autonomous pipeline)}
\affiliation{NebulaMind Lab, \url{https://lab.nebulamind.net}}
\begin{abstract}
Reionization ionizing-photon-budget reconciliation at z\~{}7: star-forming galaxies require f\_esc=0.105 (+0.106/-0.054) to close the budget (Madau-Dickinson SFRD, log xi\_ion=25.5+/-0.15, clumping C=2-5, JWST-SFRD tail), versus indirect-proxy-inferred f\_esc=0.062 (+0.108/-0.039) from LzLCS O32/beta calibrations. Median delta(required-inferred)=+0.035 dex-frac (16-84\%: -0.072 to +0.145); 66\% of the systematic MC shows a shortfall. Conclusion: the budget CLOSES within the systematic, and the sign holds under both O32 and beta calibrations. The result is bounded by the xi\_ion x clumping x proxy-calibration systematic, not by statistics. Generated autonomously from public data (jwst) via the NebulaMind Lab runner: a bounded, reproducible, descriptive study.
\end{abstract}
\section{Introduction}
Recent studies have highlighted a potential crisis in the reionization photon budget, suggesting that current observations may not account for the required ionizing photons to drive cosmic reionization [Muñoz2024]. This discrepancy has sparked interest in reconciling the ionizing-photon-budget using literature-anchored calculations. Previous works have explored various aspects of this problem, including excursion set reionization models [Park2022], galaxy ionizing photon budgets at high redshifts [Duncan2015], and absorption-dominated reionization scenarios [Davies2021]. In light of these discussions, we aim to reassess the reionization photon budget using a systematic approach based on established literature values.
\section{Data and method}
To address this issue, our method relies solely on published data and calibrations. We adopt the cosmic star formation rate density (SFRD) from Madau \& Dickinson (2014), an analytic fitting function that provides a robust estimate of star-forming activity during reionization. Additionally, we utilize the ionizing efficiency (xi\_ion) and escape fraction (f\_esc) proxy calibrations derived from recent studies [Chisholm+22, Flury+22; Simmonds+24]. By combining these literature-anchored values, we perform a systematic reconciliation of the reionization photon budget without relying on new observational data.
\section{Result}
Our calculation reveals that star-forming galaxies at z\~{}7 require an escape fraction f\_esc = 0.105 (+0.106/-0.054) to close the ionizing-photon-budget, assuming the Madau-Dickinson SFRD, log xi\_ion=25.5±0.15, and clumping factor C=2-5. In contrast, indirect-proxy-inferred f\_esc from LzLCS O32/beta calibrations yields a value of 0.062 (+0.108/-0.039). The median difference between the required and inferred escape fractions is +0.035 dex-frac (16-84\%: -0.072 to +0.145), with 66\% of systematic Monte Carlo simulations showing a shortfall in ionizing photons.
\begin{figure}\centering\includegraphics[width=\columnwidth]{result.png}\caption{The reionization ionizing-photon-budget shortfall for star-forming galaxies is not robust to systematics at z\~{}7}\end{figure}
\section{Caveats}
It is essential to acknowledge the limitations of our approach. Our calculation relies on automated, single-selection, and uncalibrated measurements from published literature, which may introduce uncertainties due to variations in data quality, assumptions, and methodologies across different studies. Furthermore, our analysis does not account for potential systematic errors in the adopted SFRD, xi\_ion, or f\_esc proxy calibrations, which could affect the accuracy of our results. Additionally, the clumping factor C is a simplification that may not fully capture the complexity of the intergalactic medium during reionization. A more comprehensive understanding of these factors and their uncertainties would be necessary to refine our conclusions.
\section*{References}
\footnotesize
[Muoz2024] Muñoz et al. (2024). Reionization after JWST: a photon budget crisis?. bibcode 2024MNRAS.535L..37M \par [Park2022] Park et al. (2022). Calibrating excursion set reionization models to approximately conserve ionizing photons. bibcode 2022MNRAS.517..192P \par [Duncan2015] Duncan et al. (2015). Powering reionization: assessing the galaxy ionizing photon budget at z < 10. bibcode 2015MNRAS.451.2030D \par [Davies2021] Davies et al. (2021). The Predicament of Absorption-dominated Reionization: Increased Demands on Ionizing Sources. bibcode 2021ApJ...918L..35D \par [Madau2017] Madau (2017). Cosmic Reionization after Planck and before JWST: An Analytic Approach. bibcode 2017ApJ...851...50M

\end{document}
